\usepackage{amsmath} % for my editor's highlighting; loaded by thmmarks option
\usepackage{amssymb}
\usepackage{tikz} % for testing tikzpicture and tikzcd
\usetikzlibrary{cd} % for testing tikzcd
\usepackage[auto-qed]{keytheorems}
\usepackage{hyperref}
\usepackage{kantlipsum} % for filler text, to check spacing
\newcommand{\kantgray}[2]{{\color{gray}\kant*[#1][#2]}}

\newkeytheoremstyle{heartstyle}{
    qed=\ensuremath{\heartsuit}
    }
\newkeytheoremstyle{clubstyle}{
    qed=\ensuremath{\clubsuit}
    }
\newkeytheoremstyle{diamondstyle}{
    qed=\ensuremath{\diamond},
    bodyfont=\normalfont
    }
\newkeytheoremstyle{bulletstyle}{
    qed=\ensuremath{\bullet},
    bodyfont=\normalfont
    }

\newkeytheorem{theorem}[
    style=heartstyle,
    parent=section,
    ]
\newkeytheorem{proposition}[
    style=clubstyle,
    parent=section
    ]
\newkeytheorem{lemma}[ % not endmarked
    parent=section
    ]
\newkeytheorem{remark}[
    style=diamondstyle,
    parent=section,
    ]
\newkeytheorem{note}[
    style=bulletstyle,
    parent=section,
    ]

% style for boring tests
\newkeytheorem{test}[
    style=diamondstyle,
    parent=section,
    ]
% maybe these definitions need to be expl3?
% e.g.,if there are unwanted space issues
\newcommand{\testtype}[2]{\texttt{#1}:% suppress extra spaces (linebreaks)
  #2}
\newcommand{\makeoutertest}[3]{
  \begin{test}
    \testtype{#1}{#2}
    #3
  \end{test}
}
\newcommand{\maketest}[3]{
  \makeoutertest{#1}{#2}{
    \begin{#1}
      #3
    \end{#1}
  }
}

%%
%% potential user commands
%%
% put thmmark at end of previous line
\NewDocumentCommand{\qedprev}{O{1}}{\par\vspace{-#1\baselineskip}\qedhere}

% like \qedprev, but without the unnecessary qedhere
\NewDocumentCommand{\qedabove}{O{1}}{\par\vspace{-#1\baselineskip}~}





\begin{document}

\section{Basic tests}
Tests of the three basic cases, as described in the package documentation.

\begin{theorem}
  Case 1: Environment ends with simple text.
\end{theorem}
\begin{theorem}
  Case 2 (unnumbered): Environment ends with display
  \[
    1 + 1 = 2.
  \]
\end{theorem}
\begin{theorem}
  Case 2 (numbered): Environment ends with equationlike
  \begin{equation}
    1 + 1 = 2.
  \end{equation}
\end{theorem}
\begin{theorem}
  Case 3 (list): Environment ends with
  \begin{itemize}
  \item a
  \item list.
  \end{itemize}
\end{theorem}
\begin{theorem}
  Case 3 (block): Environment ends with a text block.
  \begin{center}
    For example, some centered text.
  \end{center}
\end{theorem}

\section{Tests with multiple or nested environments}

These test multiple or nested boxes inside of theoremlike environments.

\begin{remark}
  Case 1: Environment ends with simple text.
  But first there is an equation
  \begin{equation}
    1+1 = 2
  \end{equation}
  and then the simple text.
\end{remark}
\begin{remark}
  Case 2 (numbered): Environment ends with display
  \begin{enumerate}
  \item but first
  \item a list
  \end{enumerate}
  \[
    1 + 1 = 2.
  \]
\end{remark}
\begin{remark}
  Case 2 (unnumbered): Environment ends with align
  following
  \begin{itemize}
  \item a
  \item list
  \end{itemize}
  and
  \[
    \texttt{ display math }
  \]
  \[
    \texttt{ again }
  \]
  \begin{align}
    x & = 2\\
    y & = 3. \quad \text{this is a very wide box that uses many words to fill the line}
  \end{align}
\end{remark}

\begin{remark}
  Case 3 (list): Environment ends with
  \begin{itemize}
  \item a
  \item list.
  \item that is nested
    \begin{enumerate}
    \item inside another
    \item list
    \end{enumerate}
  \end{itemize}
\end{remark}

\begin{note}
  Case 3 (list): Environment ends with
  \begin{itemize}
  \item a
  \item list
  \item with a display
    \[
      a^2 + b^2 = c^2
    \]
    \begin{enumerate}
    \item and there is
    \item a nested list
    \end{enumerate}
  \item but the outer list has
  \item the final item 
  \end{itemize}
\end{note}

\begin{note}\ 
  \begin{itemize}
  \item a
  \item list
    \begin{enumerate}
    \item and there is
    \item a nested list
    \item with a display
      \[
        a^2 + b^2 = c^2
      \]
    \end{enumerate}
  \end{itemize}
\end{note}
\begin{note}\ 
  \begin{itemize}
  \item a
  \item list
    \begin{enumerate}
    \item and there is
    \item a nested list
    \item with align
      \begin{align}
        a^2 + b^2 & = c^2\\
        3^2 + 4^2 & = 5^2
      \end{align}
    \end{enumerate}
  \end{itemize}
\end{note}
\begin{note}\ 
  \begin{itemize}
  \item a
  \item list
    \begin{enumerate}
    \item and there is
    \item a nested list
    \item with align
      \begin{align}
        a^2 + b^2 & = c^2\\
        3^2 + 4^2 & = 5^2
      \end{align}
    \end{enumerate}
  \end{itemize}
  Then there is some text at the end
\end{note}
\begin{note} 
  First an align
  \begin{align}
    a^2 + b^2 & = c^2\\
    3^2 + 4^2 & = 5^2
  \end{align}
  Then some text, 
  \begin{enumerate}
  \item and then
  \item a list.
  \end{enumerate}
\end{note}


\begin{note} 
  First some text,
  \begin{enumerate}
  \item and then
  \item a list.
  \end{enumerate}
  Then
  \begin{center}
    Centered text that ends with a display
    \[
      5^2 + 12^2 = 13^2.
    \]
  \end{center}
\end{note}

\subsection{some other messy tests}

\begin{theorem}
  This ends with
  \begin{enumerate}
  \item an
  \item enumerated
  \item list.
  \end{enumerate}
\end{theorem}

\begin{theorem}
  This ends with
  \begin{enumerate}
  \item an
  \item enumerated
  \item list
  \item ending in equation:
  \begin{align*}
  a&=b
  \end{align*}
  \end{enumerate}
\end{theorem}

\begin{theorem}
  This ends with
  \begin{enumerate}
  \item an
  \item enumerated
  \item list.
  \begin{enumerate}
  \item abc
  \end{enumerate}
  \item blub
  \end{enumerate}
\end{theorem}

\begin{theorem}
  This ends with
  \begin{enumerate}
  \item an
  \item enumerated
  \item list.
  \begin{enumerate}
  \item abc
  \end{enumerate}
  \end{enumerate}
\end{theorem}

\subsection{proof tests}

\begin{proof}
regular text
\end{proof}

\begin{proof}
  This ends in a display
  \[
    3^2 + 4^2 = 5^2.
  \]
\end{proof}

\begin{proof}
  This ends in a display
  \[
    3^2 + 4^2 = 5^2.
  \]
more text
\begin{align}
a&=b \\
c&=d
\end{align}
%\[a=b\]
\end{proof}

\begin{proof}
  This ends in a display
  \[
    3^2 + 4^2 = 5^2.
  \]
more text
\begin{align*}
a=b
\end{align*}
abc
\end{proof}

\begin{proof}
  This ends with
  \begin{enumerate}
  \item an
  \item enumerated
  \item list.
  \end{enumerate}
\end{proof}

\begin{proof}
  This ends with
  \begin{enumerate}
  \item an
  \item enumerated
  \item list
  \item ending in equation:
  \begin{align*}
  a&=b
  \end{align*}
  \end{enumerate}
\end{proof}

\begin{proof}
  This ends with
  \begin{enumerate}
  \item an
  \item enumerated
  \item list.
  \begin{enumerate}
  \item abc
  \end{enumerate}
  \item blub
  \end{enumerate}
\end{proof}

\begin{proof}
  This ends with
  \begin{enumerate}
  \item an
  \item enumerated
  \item list.
  \begin{enumerate}
  \item abc
  \end{enumerate}
  \end{enumerate}
\end{proof}


\section{Math environments}

Subsections are intended to group environment types according to how they are handled internally; revise as necessary.

Many of the examples and descriptions are from  the \href{https://www.ams.org/arc/tex/amsmath/amsldoc.pdf}{amsmath documentation}: Section 3, Displayed equations.

\subsection{equationlike}

These environments are full line width.
They do not use horizontal alignment markers (\&).

\maketest{equation}{
  a single equation with an automatically generated number
}{
  \int_a^b f(x)\, dx
}
\maketest{equation*}{
  unnumbered variant
}{
  \int_a^b f(x)\, dx
}
\makeoutertest{$\backslash$[..$\backslash$]}{
  The classic display math.
  Internally, amsmath redefines this as equation*
}{
  \[
    \int_a^b f(x)\, dx
  \]
}

\maketest{multline}{
  a variation of equation, for long equations with linebreaks.
  The first line is aligned left, middle lines are centered, last line is aligned right.\\
  Equation number is on the first line; end mark is on the last line.
}{
  a+b+c+d+e+f+ a+b+c+d+e+f\\
  +g+h+ g+h+ g+h+\\
  +i+j+k+l+m+n
  +i+j+k+l+m+n
}
\maketest{multline*}{
  unnumbered variant
}{
  a+b+c+d+e+f+ a+b+c+d+e+f\\
  +g+h+ g+h+ g+h+\\
  +i+j+k+l+m+n
  +i+j+k+l+m+n
}
\maketest{multline}{
  a variation of equation, for long equations with linebreaks.
  The first line is aligned left, middle lines are centered, last line is aligned right.\\
  Equation number is on the first line; end mark is on the last line.
}{
  abc\\
  dfg
}



\maketest{displaymath}{
  This is an older environment, not one of the amsmath ones; the amsmath documentation says it's functionally equivalent to equation*
}{
  \int_a^b f(x)\, dx
}

\subsection{alignlike}
These environments are full line width.
They use horizontal alignment markers (\&).

\maketest{align}{
  for two or more equations; this test has just one alignment column
}{
  \int_a^b f(x)\, dx &= 1\\
  \int_a^b g(x) + h(x)\, dx &= a+b+c
}
\maketest{align}{
  another test, with multiple alignment columns
}{
  \int_a^b f(x)\, dx &= 1
  & x &= y+z
  & p &= 0\\
  \int_a^b g(x) + h(x)\, dx &= a+b+c
  & x'+y' &= z'
  & q &= 1
}
\maketest{align*}{
  unnumbered variant; one column
}{
  \int_a^b f(x)\, dx &= 1\\
  \int_a^b g(x) + h(x)\, dx &= a+b+c
}
\maketest{align*}{
  unnumbered variant; multiple columns
}{
  \int_a^b f(x)\, dx &= 1
  & x &= y+z
  & p &= 0\\
  \int_a^b g(x) + h(x)\, dx &= a+b+c
  & x'+y' &= z'
  & q &= 1
}
\makeoutertest{alignat}{
  allows the horizontal space between equations to be explicitly specified;
  takes an argument for number of ``equation columns''.
}{
  \begin{alignat}{2}
    x& = y_1-y_2+y_3-y_5+y_8-\dots
    &\qquad& \text{by (C)}\\
     & = y'\circ y^* && \text{by (D)}\\
     & = y(0) y' && \text {by Axiom 1.}
  \end{alignat}
}
\makeoutertest{alignat*}{
  unnumbered variant
}{
  \begin{alignat*}{2}
    x& = y_1-y_2+y_3-y_5+y_8-\dots
    &\qquad& \text{by (C)}\\
     & = y'\circ y^* && \text{by (D)}\\
     & = y(0) y' && \text {by Axiom 1.}
  \end{alignat*}
}

\maketest{gather}{
  a group of numbered equations with no horizontal alignment; each line is horizontally centered.\\
}{
  \int_a^b f(x)\, dx\\
  \int_a^b f(x)g(x)h(x)\, dx
}
Note: internally, \texttt{gather} and \texttt{gather*} are handled the same way as \texttt{align}, with the other alignlike environments, even though \texttt{gather} doesn't use alignment markers (\&).
\maketest{gather*}{
  unnumbered variant
}{
  \int_a^b f(x)\, dx\\
  \int_a^b f(x)g(x)h(x)\, dx
}



\section{Lists}

\maketest{enumerate}{
  items are counted
}{
\item top
\item middle
\item bottom
}
\maketest{itemize}{
  items are not counted
}{
\item top
\item middle
\item bottom
}
\maketest{description}{
  items have labels
}{
\item[The first one:] goes here
\item[Another:] and this is the item content
\item[The final one:] these are often used for defintions
}
\makeoutertest{list}{
  This is the basic environment that others are derived from.
}{
  \begin{list}{}{}
  \item first
  \item middle
  \item last
  \end{list}
}


\section{Other blocks}

These other special cases are each handled separately.

\maketest{quote}{}{
  A quotation goes here.
  \kantgray{1}{1}
}


\maketest{center}{
  center the contents
}{
  some centered text, or equation, or other content
}
Note: the final line is off-centered, to make room for the endmark.
It is usually unnoticeable, except in cases with visually similar lines, such as the following.
\maketest{center}{
  center the contents
}{
  some centered text, or equation, or other content\\
  some centered text, or equation, or other content\\
  some centered text, or equation, or other content%
  % \newdimen{\markwidth}%
  % \settowidth{\markwidth}{ \qedsymbol}%
  % \hspace*{-2\markwidth}% not quite right
}

In cases like the above, some manual adjustment is necessary.
For example, some negative \verb|\hspace| for the width of the end mark.

Alternatively, an extra line break can move the endmark to a new line.
This is shown in the next example.
\maketest{center}{
  center the contents
}{
  some centered text, or equation, or other content\\
  some centered text, or equation, or other content\\
  some centered text, or equation, or other content\\~ % extra space necessary
}

\subsection{tabbing}

See \href{https://mirrors.ctan.org/macros/latex/base/source2e.pdf}{source2e}: File 41 ltab.dtx: Tabbing, Tabular, and Array Environments.

Explanation and examples from \url{https://latexref.xyz}.

\maketest{tabbing}{
  manually set and use tab stops
}{
  row1-col1\qquad \= row1-col2\qquad \= this one has a longer heading\qquad \= r1-c4 \\
  row2-col1 \> row2-col3  \> . \> . \\
  r3-c1 \> next field \> a \> b
}
\maketest{tabbing}{
  manually set and use tab stops
}{
  row1-col1\qquad \= row1-col2\qquad \= this one has a longer heading\qquad \= r1-c4 \\
  row2-col1 \> row2-col3  \> . \> . \\
  r3-c1 \> next field \> a \> longer entry 012
}
Above: the final longer entry pushes the end mark down.

\subsection{verbatim}

Internally, this is handled with the same code that handles lists (trivlist).
Currently this only sort of works.
The \verb|\end{verbatim}| must be separated by a space, but not a newline, from the contents.
This could have to do with the \texttt{hmode} v.s. \texttt{vmode} checks in the thmmark code.

% \begin{test} % manual, because argument expansion seems to get confused
%   \texttt{verbatim}:
%   Note that linebreaks are also treated verbatim.
% \begin{verbatim}
% Some verbatim symbols \ & $ etc. \end{verbatim}
% \end{test}

\begin{test} % manual, because argument expansion seems to get confused
  \texttt{verbatim}:
  Note that linebreaks are also treated verbatim.
\begin{verbatim}
Some verbatim symbols \ & $ etc.
More verbatim symbols \ & $ etc. More verbatim symbols \ & $ etc. More verbatim symbols \ & $ etc.
(The previous line intentionally runs off the page) \end{verbatim}
\end{test}

\begin{test} % manual, because argument expansion seems to get confused
  \texttt{verbatim}:
  (Another test)
  If the last line is \emph{just a little} too long, it will push the endmark into the margin.
\begin{verbatim}
A123456789.B123456789.C123456789.D123456789.E123456789.F123456789 \end{verbatim}
\end{test}
\begin{test} % manual, because argument expansion seems to get confused
  \texttt{verbatim}:
  (Another test)
  If the last line is \emph{a little more} too long, it will push the endmark down.
\begin{verbatim}
A123456789.B123456789.C123456789.D123456789.E123456789.F123456789. \end{verbatim}
\end{test}



\section{Unsupported}

These environments are currently unsupported.
The marks may or may not be placed as expected, depending on the surrounding environment(s).

\subsection{Unsupported math environments}



\makeoutertest{\$\$..\$\$}{
  This syntax is still used for display math.
  By default, it doesn't work well with \texttt{qedhere}; the end mark is not pushed to end of line.
}{
  $$ 
    \int_a^b f(x)\, dx
  $$
}


\maketest{eqnarray}{
  not recommended, for many reasons documented in
  \href{https://tug.org/TUGboat/tb33-1/tb103madsen.pdf}{Avoid \texttt{eqnarray}!} (TUGboat 33:1, 2012)
}{
  \int_a^b f(x)\, dx & = 1.
}


\maketest{flalign}{
  full length alignment; stretches to fill line, leaving space for equation numbers and endmark (if present)\\
  Currently unsupported.  Manual placement of via \texttt{qedhere} works fine, but automatic placement with our alignlike code causes overlap of end mark with content.
}{
  x&=y & X&=Y\\
  x'&=y' & X'&=Y'\\
  x+x'&=y+y' & X+X'&=Y+Y'% \qedhere %manual placement works correctly
}
\maketest{flalign*}{
  unnumbered variant; as with the numbered variant, manual placement treats endmark correctly, but automatic does not\\
}{
  x&=y & X&=Y\\
  x'&=y' & X'&=Y'\\
  x+x'&=y+y' & X+X'&=Y+Y' %\qedhere
}


\subsection{Math environments that go inside others}
The width of the following environments is just the width of the content.
Each of these examples is inside \texttt{equation}.
Each takes an optional argument for vertical positioning (baseline); demonstrated just with the first one.

\makeoutertest{gathered}{
  [c] (the default)
}{
  \begin{equation}
    \begin{gathered}[c]
      \int_a^b f(x)\, dx\\
      \int_a^b f(x)\, dx\\
      \int_a^b f(x)\, dx\\
      a+b+c+d = e+f+g+h+i+j+k
    \end{gathered}
  \end{equation}
}
\makeoutertest{gathered}{
  [t] (baseline top)
}{
  \begin{equation}
    \begin{gathered}[t]
      \int_a^b f(x)\, dx\\
      \int_a^b f(x)\, dx\\
      \int_a^b f(x)\, dx\\
      a+b+c+d = e+f+g+h+i+j+k
    \end{gathered}
  \end{equation}
}
\makeoutertest{gathered}{
  [b] (baseline bottom)
}{
  \begin{equation}
    \begin{gathered}[b]
      \int_a^b f(x)\, dx\\
      \int_a^b f(x)\, dx\\
      \int_a^b f(x)\, dx\\
      a+b+c+d = e+f+g+h+i+j+k
    \end{gathered}
  \end{equation}
}
\makeoutertest{aligned}{
}{
  \begin{equation}
      \begin{aligned}
      \int_a^b f(x)\, dx & = 2 \\
      \int_a^b f(x)\, dx & = x+y+z \\
      \int_a^b f(x)\, dx & = 1\\
      a+b+c+d & = e+f+g+h+i+j+k
      \end{aligned}
  \end{equation}
}
\makeoutertest{alignedat}{
}{
  \begin{equation}
    \begin{alignedat}{2}
      x& = y_1-y_2+y_3-y_5+y_8-\dots
      &\qquad& \text{by (C)}\\
       & = y'\circ y^* && \text{by (D)}\\
       & = y(0) y' && \text {by Axiom 1.}
    \end{alignedat}
  \end{equation}
}
\makeoutertest{cases}{
}{
  \begin{equation*}
    P_{r-j}=
    \begin{cases}
      0& \text{if $r-j$ is odd},\\
      r!\,(-1)^{(r-j)/2}& \text{if $r-j$ is even}.
    \end{cases}
  \end{equation*}
}

\makeoutertest{dcases}{
  This is a variant that the amsmath docs mention, but is provided by the \texttt{mathtools} package.
  It makes the case lines displaystyle instead of (the default) textstyle.
  So, maybe we don't want to support this directly (especially not if we decide not to make \texttt{cases} endmarkable in the first place).
}{
  \break [\emph{no example here}]
}


% \subsection{inner math environments}
% These environments go inside other displayed math environments.\\
% They're not directly supported, because their end marks can be handled by the surrounding environments.

\makeoutertest{split}{
  for a \emph{single} equation split among multiple lines, like \texttt{multline}, but with some alignment points.\\
  This \texttt{split} is inside \texttt{equation}.
  Note: the end mark is missing.
}{
  \begin{equation}
    \begin{split}
      H_c
      &=\frac{1}{2n} \sum^n_{l=0}(-1)^{l}(n-{l})^{p-2}
        \sum_{l _1+\dots+ l _p=l}\prod^p_{i=1} \binom{n_i}{l _i}\\
      &\quad\cdot[(n-l )-(n_i-l _i)]^{n_i-l _i}\cdot
        \Bigl[(n-l )^2-\sum^p_{j=1}(n_i-l _i)^2\Bigr].
    \end{split}
  \end{equation}
}

\makeoutertest{split}{
  This \texttt{split} is inside \texttt{gather}, and the endmark is present.
}{
  \begin{gather}
    % a + b = c\\
    \begin{split}
      H_c
      &=\frac{1}{2n} \sum^n_{l=0}(-1)^{l}(n-{l})^{p-2}
        \sum_{l _1+\dots+ l _p=l}\prod^p_{i=1} \binom{n_i}{l _i}\\
      &\quad\cdot[(n-l )-(n_i-l _i)]^{n_i-l _i}\cdot
        \Bigl[(n-l )^2-\sum^p_{j=1}(n_i-l _i)^2\Bigr].
    \end{split}%\\
    % \int_a^b f(x)\, dx
  \end{gather}
}
Note: The endmark is present in one \texttt{split} example, but not the other.
This seems like a bug.


\makeoutertest{pmatrix}{
  other matrix environments are likewise not (directly) supported: \texttt{smallmatrix}, \texttt{bmatrix}, \texttt{Bmatrix}, \texttt{vmatrix}, and \texttt{Vmatrix}.
  Just let the surrounding environment handle the end mark.
}{
  \[\begin{pmatrix}
    a&b\\ c&d
  \end{pmatrix}\]
}

\makeoutertest{array}{
  works inside displayed math; this one is in equation*.
  Just let the surrounding environment handle the end mark.
}{
\begin{equation*}
  \chi(x) =
  \left|              % vertical bar fence
    \begin{array}{ccc}
      x-a  &-b  &-c  \\
      -d   &x-e &-f  \\
      -g   &-h  &x-i
    \end{array}
 \right|
\end{equation*}
}






\subsection*{Others}

A few more examples that are often wrapped in other environments such as \texttt{center}.


\subsubsection*{tabular}

\makeoutertest{tabular}{
  a table; wrap in another environment to position endmark
}{
\begin{tabular}{l|l}
  \textit{Player name}  &\textit{Career home runs}  \\
  \hline
  Hank Aaron  &755 \\
  Babe Ruth   &714 \\
  Hank Aaron (repeat)  &755 \\
  Babe Ruth (repeat)  &714
\end{tabular}
}
\begin{test}
  \texttt{tabular}: 
  a table in a \texttt{center}, with experimental command \verb|\qedabove| to manually position the endmark
  \begin{center}
    \begin{tabular}{l|l}
      \textit{Player name}  &\textit{Career home runs}  \\
      \hline
      Hank Aaron  &755 \\
      Babe Ruth   &714 \\
      Hank Aaron (repeat)  &755 \\
      Babe Ruth (repeat)  &714
    \end{tabular}
  \end{center}\qedabove[1.3]
\end{test}

Note: \verb|\qedabove| can mess up spacing of the following lines, such as this one.
\kantgray{1}{2}



\subsubsection*{tikz environments}
If these are wrapped in other environments like \texttt{equation}, then they won't need direct support

\makeoutertest{tikzpicture}{
  often used for diagrams, but of course requires tikz\\
}{
  \hspace*{.7\textwidth}
  \begin{tikzpicture}
    \draw (0,0) -- ++(1,1) -- ++(1,-1);
  \end{tikzpicture}
}

% doesn't work with our test macros
\begin{test}
  \texttt{tikzcd}:
  another popular one, also requires extra packages\\
  This is a \texttt{tikzcd} inside an \texttt{equation}.
  \begin{equation}
    \begin{tikzcd}
      A \arrow[r, "\phi"] \arrow[d]
      & B \arrow[d, "\psi"] \\
      C \arrow[r, "\eta"]
      & D
    \end{tikzcd} 
  \end{equation}
\end{test}


\section{Tests of nested end marked environments}
Nesting end marked environments inside eachother is not officially supported, but sort of works.
End marks for successive closed environments are placed on new lines.

\begin{theorem}
  regular text
  \begin{lemma}
    blub; lemma is not endmarked
  \end{lemma}
  \begin{lemma}
    abc
    \begin{equation*}
      abc
    \end{equation*}
    \begin{enumerate}
    \item abc
      \begin{itemize}
      \item blub
      \end{itemize}
    \item jhdgd
    \end{enumerate}
  \end{lemma}
\end{theorem}

Text following theorem environment.
\kantgray{2}{1-2}

\begin{theorem}
  regular text
  \begin{remark}
    blub; remark \emph{is} endmarked
  \end{remark}
\end{theorem}

Text following theorem environment.
\kantgray{3}{1-2}

\begin{theorem}
  regular text
  \begin{remark}
    blub; remark \emph{is} endmarked
  \end{remark}
  \begin{lemma}
    abc; lemma is not endmarked
    \begin{equation*}
      abc
    \end{equation*}
    \begin{enumerate}
    \item abc
      \begin{itemize}
      \item blub
      \end{itemize}
    \item jhdgd
    \end{enumerate}
  \end{lemma}
\end{theorem}


Text following theorem environment.
\kantgray{6}{1-2}

 
Note: When a proof is nested inside another environment, as in the following test, its qed end mark changes to that of the containing environment.
Here, environments are nested inside proofs nested inside other environments.

\begin{theorem}
  \kantgray{1}{1}\\
  This proof is inside the theorem environment.
  \begin{proof}
    \kantgray{2}{3-5}
  \end{proof}
\end{theorem}



\begin{theorem}
  \kantgray{1}{1}\\
  This proof is inside the theorem environment.
  \begin{proof}
    \kantgray{3}{1}
    \begin{proposition}
      This proposition is inside the theorem proof.
      \kantgray{5}{3}
      \begin{proof}
        This proposition proof is inside the proposition environment.
        \kantgray{5}{4-5}
        [End proposition proof; end proposition.]
      \end{proof}
    \end{proposition}
    [Theorem proof continues]
    \kantgray{6}{2}
    \begin{remark}
      Remark still inside theorem proof.
      \kantgray{6}{3}
      [End remark; end theorem proof; end theorem.]
    \end{remark}
  \end{proof}
\end{theorem}

Paragraph after theorem environment.
\kantgray{2}{3-4}

\clearpage
\section{Other Tests}

Here are experimental tests of:
\begin{enumerate}
\item ways to disable the thmmark for one box
\item ways to put the thmmark on a previous line
\item thmmarks behavior when content is full text width
\end{enumerate}
These are buggy and incomplete.

\newcommand{\noqed}{\renewcommand{\qedsymbol}{}}

\begin{test} 
  Use \verb|\par| and negative \verb|\vspace| to manually set endmark at bottom of environment.
  Note: This will cause overlap with content, if the previous line is too long, so it should only be placed manually.
  \begin{equation*}
    \chi(x) =
    \left|              % vertical bar fence
      \begin{array}{ccc}
        x-a  &-b  &-c  \\
        -d   &x-e &-f  \\
        -g   &-h  &x-i
      \end{array}
    \right|% \noqed 
  \end{equation*}
  \par\vspace{-1.3\baselineskip}~
\end{test}

Paragraph after environment, to check spacing.
\kantgray{2}{1-2}

\begin{test} (duplicate, with no extra commands, to check spacing)
  \begin{equation*}
    \chi(x) =
    \left|              % vertical bar fence
      \begin{array}{ccc}
        x-a  &-b  &-c  \\
        -d   &x-e &-f  \\
        -g   &-h  &x-i
      \end{array}
    \right|
  \end{equation*}
\end{test}

Paragraph after environment, to check spacing.
\kantgray{2}{1-2}


\begin{test} 
  \verb|\noqed| does not work here, at all
  \begin{align}
    x & = 2\\
    y & = 3\\
    P_{r-j} & =
    \begin{cases}
      0& \text{if $r-j$ is odd},\\
      r!\,(-1)^{(r-j)/2}& \text{if $r-j$ is even}\\
      \text{other} & \text{case, for testing.}\noqed
    \end{cases} \noqed
  \end{align}\noqed
  a line of text \noqed 
\end{test}

\begin{test} 
  \verb|\noqed| does not work here
  \begin{align}
    x & = 2\\
    y & = 3\noqed
  \end{align}
\end{test}

\subsection{tikzpicture}

Some tests of \texttt{tikzpicture} in equations.
This is how I (NGJ) often make numbered commutative diagrams.


\begin{test}
  \texttt{tikzpicture} in \texttt{equation}\\
  baseline (0,0)
  \begin{equation}
    \hspace*{.7\textwidth}
    \begin{tikzpicture}
      \draw (0,0) -- ++(1,1) -- ++(1,-1);
    \end{tikzpicture}
  \end{equation}
\end{test}
Spacing test.
\kantgray{3}{1-2}
\begin{test}
  another \texttt{tikzpicture} in \texttt{equation}\\
  baseline (1,1)
  \begin{equation}
    \hspace*{.7\textwidth}
    \begin{tikzpicture}[baseline={(1,1)}]
      \draw (0,0) -- ++(1,1) -- ++(1,-1);
    \end{tikzpicture}
  \end{equation}
\end{test}
Spacing test.
\kantgray{3}{1-2}
\begin{test}
  another \texttt{tikzpicture} in \texttt{equation}\\
  baseline (1,1); use \verb|\par| and negative \verb|\vspace| to reposition endmark
  \begin{equation}
    \hspace*{.7\textwidth}
    \begin{tikzpicture}[baseline={(1,1)}]
      \draw (0,0) -- ++(1,1) -- ++(1,-1);
    \end{tikzpicture}
  \end{equation}\qedprev[1.3]
\end{test}
Spacing test.
\kantgray{3}{1-2}

\subsection{full-width equations}

\begin{test}
  This is a test of an equation that is almost the full text width.
  The equation number appears on the line above the long equation; that seems good.
  The endmark also appears on that line above; that seems bad.
  For a manual fix: add a space like ``\verb|\ |'' (backslash space) to move the end mark to a new blank line.
  \begin{equation}
    \sum_i a_i + \sum_i b_i
    + \sum_i a_i + \sum_i b_i
    + \sum_i a_i + \sum_i b_i
    + \sum_i a_i + \sum_i b_i
    + \sum_i a_i + \sum_i b_i
    + c + d + e
  \end{equation}%\ % add new line; endmark appears on new line
\end{test}
\begin{test}
  Same test, ending with ``\verb|\ |'' (backslash space).
  \begin{equation}
    \sum_i a_i + \sum_i b_i
    + \sum_i a_i + \sum_i b_i
    + \sum_i a_i + \sum_i b_i
    + \sum_i a_i + \sum_i b_i
    + \sum_i a_i + \sum_i b_i
    + c + d + e
  \end{equation}\ % add new line; endmark appears on new line
\end{test}

\begin{test}
  Same test, with \texttt{align}.
  Note: spacing between $+$ signs here is different from \texttt{equation}.
  End mark is printed over end of second line
  \begin{align}
    \sum_i a_i + \sum_i b_i
    + \sum_i a_i + \sum_i b_i
    + \sum_i a_i + \sum_i b_i
    + \sum_i a_i + \sum_i b_i
    % + \sum_i a_i + \sum_i b_i
    & + c + d + e + f \\
    & = 0123456789.01x
  \end{align}%\ % add new line; endmark appears on new line
\end{test}
\begin{test}
  Similar test, with \texttt{align}.
  In this one, the second equation number is also moved up, and the end mark is printed on that line.
  \begin{align}
    \sum_i a_i + \sum_i b_i
    + \sum_i a_i + \sum_i b_i
    + \sum_i a_i + \sum_i b_i
    + \sum_i a_i + \sum_i b_i
    % + \sum_i a_i + \sum_i b_i
    & + c + d + e + f\\
    \sum_i a_i + \sum_i b_i
    + \sum_i a_i + \sum_i b_i
    + \sum_i a_i + \sum_i b_i
    + \sum_i a_i + \sum_i b_i
    % + \sum_i a_i + \sum_i b_i
    & + c + d + e + f
  \end{align}%\ % add new line; endmark appears on new line
\end{test}

\clearpage
\section{Tall lines in align and friends}

These tests demonstrate end mark placement when an \texttt{align} or similar environment has a very tall line.
By default, the end mark is placed on the baseline, aligned with an equation number, if present.
Some of these tests experiment with different ways to manually adjust the vertical alignment.

\maketest{align}{
  with \texttt{cases}
}{
  \text{words to push this content to the right edge}\qquad\qquad x & = 2\\
  y & =
  \begin{cases}
    \int_a^b f(x)\,dx \\
    \int_a^b f(x)\,dx \\
    \int_a^b f(x)\,dx \\
    \int_a^b f(x)\,dx 
  \end{cases}
}
Spacing test.
\kantgray{3}{1-2}
\begin{test} 
  Similar to above, but manually reposition endmark with\\
  \verb|  \par\vspace{-1.6\baselineskip}~|\\
  after \verb|\end{align}|
  \begin{align}
    \text{words to push this content to the right edge}\qquad\qquad x & = 2\\
    y & =
        \begin{cases}
          \int_a^b f(x)\,dx \\
          \int_a^b f(x)\,dx \\
          \int_a^b f(x)\,dx \\
          \int_a^b f(x)\,dx 
        \end{cases}
  \end{align}\qedprev[1.6]\hspace{0.001pt}\par\vspace{.6\baselineskip}
  % note: \par above and positive \vspace seem to be necessary
  % but then next line is indented because of \par
  What if there is another line of text?  Is the endmark placed above?
\end{test}

Spacing test.
\kantgray{3}{1-2}

Notes for the previous example.
\begin{enumerate}
\item Ending with \verb|\qedabove| also works.
\item Ending with \verb|\qedhere| doesn't work; the endmark is on the baseline.
\item Ending with \verb|\qedhere| (as in \verb|\qedprev|) followed by \verb|\hspace| \emph{does} work; the endmark appears after the fourth \texttt{cases} row.\\
  It doesn't matter how much \texttt{hspace}, as long as it's a nonzero amount.
  This is probably something to do with \texttt{vmode} v.s.~\texttt{hmode}, and if so then there is probably a better way to handle it.
  That's beyond what I currently understand though.
\item Spacing after the environment doesn't seem to be affected by the negative \texttt{vspace}, but lines \emph{inside} the environment are affected.
\item Maybe the whole negative \texttt{vspace} idea is just not great; is there a better alternative?
\end{enumerate}

\end{document}
